\chapter{选题依据}

\section{研究背景}
\placeholder{从国家粮食安全、智慧农业/数字农业、规模化玉米生产与苗期精准监测需求切入。建议先提出万亩级生产场景的管理问题，再说明现有人工调查或局部低空监测难以同时满足“精细度、覆盖范围与效率”的矛盾。}

\section{研究目的与意义}
\subsection{研究目的}
\placeholder{围绕万亩级玉米苗期监测，说明拟建立“低空逐株—高空结构—单株表型—跨域应用”的多尺度无人机遥感方法体系。}

\subsection{理论意义}
\placeholder{说明多尺度视觉表征、小目标感知、行结构建模、跨区域泛化等方面的科学或方法学意义。}

\subsection{应用意义}
\placeholder{说明对出苗数量核查、缺苗断垄定位、苗情空间差异分析、田间管理决策与规模化生产管理的作用。}

\section{国内外研究现状}
\subsection{无人机低空玉米幼苗目标检测与出苗计数}
\placeholder{综述 12--20 m 等低空尺度下的小目标检测、密集目标计数、超高分辨率影像切片与尺度适配研究。}

\subsection{中高空作物行、断行与缺苗区域识别}
\placeholder{综述 25--50 m 尺度下作物行提取、行结构约束、断行/缺苗区域识别及大区域快速普查。}

\subsection{单株冠层实例分割与影像表型测量}
\placeholder{综述单株实例分割、冠层面积/形态/覆盖特征提取，以及仅依赖影像可稳定获得的表型指标。}

\subsection{跨区域、跨种植制度与跨域泛化}
\placeholder{综述域适应、域泛化、跨地区/跨季节/跨品种迁移，以及春玉米到夏玉米等场景变化。}

\section{文献评价与问题凝练}
\placeholder{不要只做“研究较少”的泛化结论。建议归纳为 3--4 个可验证的问题，例如：低空方法难以直接扩展到大区域；高空影像中单株目标退化而行结构仍可利用；表型特征受分辨率与遮挡影响；模型跨区域稳定性不足。}

\section{本研究的切入点}
\placeholder{由上述问题自然推导到四项研究任务，并明确四项任务不是孤立模块，而是从精细监测到规模化应用、再到泛化验证的递进关系。}
